\section{The Perennial Wisdom}

Sanatana has a dual meaning: perennial in the sense of perduring over time, but also ``primordial'', since what perdures had to exist at the beginning. Specifically, it lasts through a cycle, i.e., a Manvantara which is the duration of manifestation of a terrestrial humanity. Clearly, this is not alien to the Western tradition which understands that duration as starting with Adam and ending in the final judgment. Perennial is not the same as ``eternal'' since the latter means ``beyond time''. Nevertheless, Guenon understands Sanatana Dharma as the ``reflection'' in our world of the eternal.

Guenon next questions whether the Sanatana Dharma is equivalent to the Philosophia Perennis as understood by the Medievals. While Sanatana and Perennis are equivalent, Guenon claims that Philosophia is not exactly the same as dharma. The latter term refers to doctrine that is not understood as a theory, but ``as a knowledge which must be effectively realized'' and its applications extend to all modalities of human life, without exception.

\begin{quotex}
This cannot be passed over, since the modern mind rejects this, even those who claim to adhere to, or to be ``influenced'' by, tradition. In a traditional society, there is no distinction between the sacred and the profane. A ``religion'' is not restricted to the private sphere nor, a fortiori, is it a matter of choice or of conscience for the individual. Religion ties the community to the eternal and to reject or deny the religion is to ex-communicate oneself, i.e., to put oneself outside the community. All activities are regulated by religious law and ritual, things are clean or unclean, taboo or allowed. These are described in Fustel de Coulanges' \emph{Ancient City}. When the neo-pagans start making decisions based on lightning strikes, the flights of birds, or the entrails of chickens, then they would show some understanding. But first of all, they must keep the hearth burning in their houses in memory of their ancestors. Unfortunately, they are embarrassed by their ancestors, just as the moderns are.
\end{quotex}


Dharma, moreover, implies the idea of permanence and stability of beings; specifically, regarding the preservation of beings, dharma resides in the conformity with their essential nature. This is certainly a key notion in Scholastic metaphysics which was concerned with understanding the essential nature of things. Moreover, a certain people or period also has its dharma; this notion was held by the Medievals who expressed typically in terms of angels or patron saints for nations, regions, and cities. Ultimately, there is a dharma for the human race, i.e., human nature, again an idea found among the Scholastics.

Dharma is the law proper to a cycle, formulated by the cosmic intelligence that reflects the Divine Will and expresses universal Order in it. In the West, this was called the Logos, the intelligence through which all things were made. Hence, Dharma is conformity to this cosmic order

After expanding on this point, Guenon equates the Sanatana Dharma to the Primordial Tradition. Through decline, this tradition has become hidden and inaccessible to ordinary humanity. It is the ``first source and the common fund of all particular traditional forms, which proceed from it by adaptation to special conditions or a people or an epoch.'' However, no particular tradition form is identical with the Primordial Tradition but is rather a veiled expression of it.

Hence, and this important fact is too often ignored, the Sanatana Dharma is equivalent to no particular traditional form, including specifically the Hindu tradition. Is the problem of the white Hindu priest clear now? First of all, he rejects and condemns the particular traditional form that is the Dharma for his people in their geographical location at this point in time. Furthermore, he assumes the Sanatana Dharma is the same as Hinduism. At best, this is a relative truth, because every other traditional form could rightly make the same claim. In effect, any of the forms could serve to express the perennial wisdom; Guenon chose the Vedas, whereas we prefer to use the Western tradition. What is dharma to one man at a particular place is adharma to a different man in a different place.

Guenon then equivocates on this point and seems to deviate from his usual cold Gallic logic. He asserts, despite what he just indicated, that the Sanatana Dharma \emph{appears} to be bound more closely to the Hindu tradition, since it derives most directly from the Primordial Tradition and is a ``sort'' of external continuation of it. ``Appear'' and ``sort'' are weak words that indicate nothing more than opinion. Certainly, appearances are deceiving, especially in this case, and results perhaps from merely linguistic conditions. Let is call it instead the Perennial Wisdom.

First of all, as we indicated at the beginning, the Perennial Wisdom is the reflection in time of the eternal; i.e., it ``derives'' from the above, not from the before. Otherwise, traditions that lack a direct historical connection to Hinduism could never embody that Wisdom. Secondly, there is no such continuity, as the fall into the Kali Yuga was a definitive rupture with what came before. That is why there is no historical record or scientific proof of a superior age prior to the Kali Yuga.

Most importantly, it is false to assume that what is closest in time is necessarily closer in truth. Just the opposite, in fact. The definition of the Kali Yuga is the time in which the Primordial Tradition was forgotten. That forgetfulness is immediate, for to ``remember'' the Primordial Tradition is the same as to ``know'' it. You may have forgotten your wife's anniversary, but as soon as you remember it, you simultaneously know it. The Western wisdom says that the owl of Minerva flies at dusk, so that a more complete understanding comes only after things can be seen in their full perspective.

Conversely, being distant in time does not necessarily mean more forgetfulness. When my mother was in the last stages of Alzheimer, she could not remember what she had for breakfast. However, I would sing old songs from the 30s and 40s with her. Those she would remember and she would sing along with me.


\flrightit{Posted on 2013-10-31 by Cologero}

\begin{center}* * *\end{center}

\begin{footnotesize}
\begin{sffamily}
\texttt{kenneth shaw on 2013-10-31 at 09:38 said:}

My sense of things is that, In the West, ebergy of the Tradition is transmitted through the image of Jesus. He functions as our Paramapurusa and Paramesvara.

Just recently I noticed that the name ``Jesus'' is linguisticly and structuraly indentical to Isha... which is a shortened form of Ishvara.

\hfill

\texttt{Logres on 2013-10-31 at 23:01 said:}

Oh, I finally see what you mean by ``remember'' : you don't mean some experiential ``flashback'' of a past life. You mean that to truly know something inwardly, you must be remembering it to begin with.

\hfill

\texttt{scardanelli on 2013-11-01 at 10:10 said:}

Logres,

I tend to see it as Tomberg describes it, as consisting of both a vertical and horizontal memory. These are the two dimensions in which we can know- both symbolically or spatially and mythologically or related to time. To know an event from the past (or the future for that matter... see Cologero's post on the The Future of Prophecy) is to remember or know it in a horizontal fashion, and to know that which is above is to remember or know in a vertical fashion. We are bound by time and space, so these factors restrict our way of knowing. Yet the spirit is not bound by such things, so it is not necessarily remembering a past life, but a knowing that is not restricted by time and space. At least, that is my limited understanding.

\hfill

\texttt{Space Lizard on 2014-12-26 at 11:32 said:}

``Is the problem of the white Hindu priest clear now?''

What about the white Sufi?


\hfill

\end{sffamily}
\end{footnotesize}

